\title[Do trained models keep attractors hidden?]{Do trained models keep attractors hidden? Forced spurious equilibria and a survival law for hiddenness}

\author{Maksim Elisov}
\email{maksimelisov2003@gmail.com}
\affiliation{Samara National Research University, Moskovskoye Shosse 34, Samara 443086, Russia}

\date{\today}

\begin{abstract}
An attractor is hidden if its basin avoids small neighborhoods of all equilibria, so it cannot be found from the
equilibria. We ask whether a trained model that reproduces a hidden attractor keeps it hidden from its own
equilibria; mostly it does not. Across 3710 hidden attractors reproduced by models of four classes trained on
six regimes, the share still hidden falls from 0.99 for perceptrons without unstable equilibria in the region of the true
system to about one third with one and to 0.03 or less with nine or more; model class, reservoir size and depth act mainly
through this number. Degree theory fixes the bookkeeping: a bounded model with nondegenerate equilibria has an odd
number of them, and every extra pair contains an unstable one. If the
index sum $s$ of the true equilibria is not one, as for Chua's circuit ($s=-1$), a bounded model must miss a true
equilibrium or add spurious ones; with no true equilibria ($s=0$) it always has one, and the reproduced attractor is
self-excited in all 29 cases. An equilibrium excites an attractor exactly when its local unstable manifold reaches the
closure of the basin. Transients toward the visible attractor raise survival, and a training penalty on unstable
equilibria raises it in two-layer perceptrons from 0.27 to 0.994 at unchanged median one-step error. Survival is not
fidelity: few hidden reproductions also have the visible equilibrium and a correct basin. Generic models do not infer
hidden attractors the data do not touch.
\end{abstract}

\maketitle

\begin{quotation}
Some attractors of nonlinear systems are hidden: no trajectory started near an equilibrium leads to them, so the usual
analysis, which starts at the equilibria, misses them. Neural networks trained on time series are now widely used in place
of equations, and they can reproduce such an attractor. This paper asks whether the attractor also stays hidden inside the
model. Often it does not. A trained model has equilibria of its own, often more than the true system, and in some systems
topology forces extra ones on any model with bounded output that copies the true equilibria. An attractor becomes
reachable from an equilibrium as soon as one of the unstable directions of that equilibrium leads into its basin. Across
ten thousand models of four kinds, the chance that the attractor stays hidden is set mainly by a single number, the count
of unstable equilibria in the region where the system lives. Generic models also do not discover hidden attractors that their data never
approach.
\end{quotation}
